% Options for packages loaded elsewhere
\PassOptionsToPackage{unicode}{hyperref}
\PassOptionsToPackage{hyphens}{url}
\PassOptionsToPackage{dvipsnames,svgnames,x11names}{xcolor}
%
\documentclass[
  11pt,
  a4paper,
]{ltjsarticle}
\usepackage{amsmath,amssymb}
\usepackage{iftex}
\ifPDFTeX
  \usepackage[T1]{fontenc}
  \usepackage[utf8]{inputenc}
  \usepackage{textcomp} % provide euro and other symbols
\else % if luatex or xetex
  \usepackage{unicode-math} % this also loads fontspec
  \defaultfontfeatures{Scale=MatchLowercase}
  \defaultfontfeatures[\rmfamily]{Ligatures=TeX,Scale=1}
\fi
\usepackage{lmodern}
\ifPDFTeX\else
  % xetex/luatex font selection
\fi
% Use upquote if available, for straight quotes in verbatim environments
\IfFileExists{upquote.sty}{\usepackage{upquote}}{}
\IfFileExists{microtype.sty}{% use microtype if available
  \usepackage[]{microtype}
  \UseMicrotypeSet[protrusion]{basicmath} % disable protrusion for tt fonts
}{}
\makeatletter
\@ifundefined{KOMAClassName}{% if non-KOMA class
  \IfFileExists{parskip.sty}{%
    \usepackage{parskip}
  }{% else
    \setlength{\parindent}{0pt}
    \setlength{\parskip}{6pt plus 2pt minus 1pt}}
}{% if KOMA class
  \KOMAoptions{parskip=half}}
\makeatother
\usepackage{xcolor}
\usepackage[margin=24mm]{geometry}
\setlength{\emergencystretch}{3em} % prevent overfull lines
\providecommand{\tightlist}{%
  \setlength{\itemsep}{0pt}\setlength{\parskip}{0pt}}
\setcounter{secnumdepth}{5}
\ifLuaTeX
  \usepackage{selnolig}  % disable illegal ligatures
\fi
\IfFileExists{bookmark.sty}{\usepackage{bookmark}}{\usepackage{hyperref}}
\IfFileExists{xurl.sty}{\usepackage{xurl}}{} % add URL line breaks if available
\urlstyle{same}
\hypersetup{
  colorlinks=true,
  linkcolor={blue},
  filecolor={Maroon},
  citecolor={Blue},
  urlcolor={blue},
  pdfcreator={LaTeX via pandoc}}

\author{}
\date{}

\begin{document}

\hypertarget{ux7ddaux5f62ux5199ux50cf}{%
\section{線形写像}\label{ux7ddaux5f62ux5199ux50cf}}

写像 \texttt{T:V→W} が線形であるとは

\[
T(x+y)=T(x)+T(y),\qquad T(\alpha x)=\alpha T(x)
\]

を満たすことです。

まとめれば

\[
T(\alpha x+\beta y)=\alpha T(x)+\beta T(y)
\]

です。

\hypertarget{ux306aux305cux3053ux306eux6761ux4ef6ux304cux91cdux8981ux304b}{%
\subsection{なぜこの条件が重要か}\label{ux306aux305cux3053ux306eux6761ux4ef6ux304cux91cdux8981ux304b}}

線形写像では、基底ベクトルがどこへ移るかだけ分かれば、すべてのベクトルの行き先が決まります。

\[
x=\sum_i c_i b_i
\]

なら

\[
T(x)=\sum_i c_iT(b_i)
\]

だからです。

これが行列で線形写像を有限個の数として記述できる理由です。

\hypertarget{ux884cux5217ux306fux57faux5e95ux4f9dux5b58}{%
\subsection{行列は基底依存}\label{ux884cux5217ux306fux57faux5e95ux4f9dux5b58}}

同じ \texttt{T}
でも基底が変われば表現行列は変わります。しかし写像そのもののランクや固有値など、基底変換に対して本質的に保存される量があります。

\end{document}
